%!TEX root = main.tex
%DECLARATION OF SOME CONSTANTS
\newcommand{\nodesNumber}{500}
\newcommand{\densitiesNumber}{8}
\newcommand{\nrMessages}{100}
%END---

\section{Experimental Scenario}
\label{sec:scenarios}

Our experiments are inspired by the scenario of a crowd of humans carrying portable communication devices (e.g., phones) in an area with no infrastructure connectivity.
The dissemination of information from one individual to all the others is critical, and uses ad-hoc connections between devices.
%
Our evaluation of broadcast algorithms for MANETs is mainly driven by two factors: \emph{network density} and \emph{nodes mobility}.
%
We describe in the following the parameters used in our study.

\smallskip
\noindent
\textbf{Simulator.}
We rely on OMNET++4.6/INET3.3, a simulation framework for the evaluation of networking algorithms~\cite{Varga:2008:OOS:1416222.1416290}.
%
This tool provides a full implementation of the IEEE 802.11g standard.
%
In addition, it contains models for the mobility of nodes and for energy consumption.

\smallskip
\noindent
\textbf{Communication Scenario.}
%\subsection*{Communication Scenario}
Simulated participants use their devices to broadcast messages without external coordination and may attempt to do so at nearly the same time.
%
We set our experiments to perform \nrMessages{} broadcast messages.
The frequency between broadcasts is $0.5$s.
%
The source of each message is randomly selected.
%
The content of each message is randomly generated, unique and with a length of 32~Bytes. 
%
%These two values allow us to evaluate the protocols under heavy and light network load.
%

\smallskip
\noindent
\textbf{Nodes Density.}
%\subsection*{Nodes Density}
We are interested in generating scenarios with different nodes densities.
%
%In particular, the density in the vicinity of each node must equal a given value plus/minus a threshold.
%
We consider three parameters: the transmission range (Tx), the map size and the number of nodes.
%
As shown in~\cite{Penrose2003}, the dependency between these parameters allows us to build networks with a given density, by using random geometric graphs.
%
In our experiments, all nodes have a Tx of 20~m.\footnote{Theoretically, IEEE 802.11 allows communication in a range between 20 and 140 meters, depending on the line of sight.
Previous empirical tests show that it can be between 20 and 30 m in indoor or crowded places~\cite{Garbinato2010}.}
The number of nodes $N$ is fixed to \nodesNumber{}. % , which is the mean value reported by previous studies on the size of manifestations in Western Europe~\cite{kriesi1995new}. %% ER useless.
%
To create a network with a given density, we set Tx and $N$ and compute the map size.
%
Then, we generate a laying on the map by using random geometric graphs.
%
We instantiate \densitiesNumber{} networks of \nodesNumber{} nodes each and set the densities to  {5,10, \dots, 40} nodes, a range of values that is frequently found in the literature.
%
This leads to a map size ranging from $125 \times 125$m to $354 \times 354$m.

\smallskip
\noindent
\textbf{Mobility.}
%\subsection*{Mobility}
We use the \emph{truncated-levy walk} (LW) mobility model, as it characterizes well the mobility traces of humans carrying wireless devices~\cite{Rhee:2011}.
%
This model generates traces composed of frequent short walks and occasional rides to distant locations.
%
We define two scenarios: One without mobility (\emph{static}) and another with mobility (\emph{dynamic}) at a walking speed of $1.4$ $m/s$ -- based on an estimation of the preferred walking speed~\cite{Browning390}.
%
We generate the dynamic scenario starting from the static one and applying the mobility based on the model.
We check that the network remains connected for the duration of the experiment, or discard the data and regenerate using a new random seed.

%
% The generation of \emph{static} scenarios was presented in the previous subsection.
% %
% To generate \emph{dynamic} scenarios, we begin with a \emph{static one}.
% %
% Then, we use the LW model to generate mobility traces for each node.
% %
\deleted{
The parameters of the LW model are set in such a way that the average node speed is close to .
}
%
\deleted{To do so, we use the following values to feed the mobility model:}

\deleted{
\begin{table}[ht]
	\footnotesize
	\centering

	\begin{tabular}{lcp{3.5cm}}
		\textbf{Parameter} & \textbf{Value} & \textbf{Description} \\

		\midrule
		FL\_EXP & -3.9 & Exponent of the flight length distribution \\
		FL\_MAX & 50 & Maximum value of the flight length distribution \\
		WT\_EXP & -1.8 & Exponent of the waiting time distribution \\
		WT\_MAX & 100 & Maximum value of the waiting time distribution \\
	\end{tabular}

	\caption{Configuration of the LW mobility model.}
	\label{tbl:mobility-parameters}
\end{table}
}

\deleted{
Afterwards, we perform several checks on the traces to guarantee that: i) the average speeds are indeed close to $1.4$ $m/s$, ii) the node density is the desired one during all the simulation time, and iii) the topology remains connected all the time.~\footnote{Except for a very low density (5) where we are unable to generate a topology that remained connected during all the experiments.}
%
Should any of these checks fail, we discard the mobility traces.
}

\smallskip
\noindent
\textbf{Energy Settings.}
This metric depends on the wireless communication chip used.
We model power consumption using the specifications of a commercial chip found in modern mobile devices.
%
We use Broadcom BCM4329 chip~\cite{BCM4329}, which supports Wi-Fi 802.11g on the 2.4 GHz frequency band.
%
Its energy consumption per operation modes is as follows: sleep (648~nW), listen (244.8~mW), receive (295.2~mW), transmit (1206~mW).
%
These values are used for the \emph{StateBasedEnergyConsumer} consumption model integrated with INET.

\deleted{
\begin{table}[ht]
	\footnotesize
	\centering
	
	\begin{tabular}{lr}
		\textbf{Mode} & \textbf{Consumption} \\
		
		\midrule
		%Off	  & nW & 57.6 \\
		Sleep  & 648 nW \\
		Receiver Listen & 244.8 mW \\
		Receiver Active (RX) & 295.2 mW \\
		Transmitter Active (TX) & 1206 mW \\
	\end{tabular}
	
	\caption{Power consumption per mode of the Broadcom BCM4329 Wi-Fi chip.}
	\label{tbl:wifi-chip}
\end{table}

To simulate this chip, we use the \emph{StateBasedEnergyConsumer} model provided by INET.
%
This model measures how many Joules were consumed based on the time the chip spends on each mode (as described in Table~\ref{tbl:wifi-chip}).
}
%
\deleted{Formally, this energy consumption model uses three functions: $c : M \mapsto \mathbb{R} $, $r_m: \mathbb{R} \mapsto M$ and $t_m: \mathbb{R} \mapsto M$ where $M$ is the set of operation modes of the chip, $c$ is the consumption of the chip on each mode, and functions $r_m, t_m$ map from simulation time to operation modes of the \emph{receiver} and \emph{transceiver}, respectively.
%
INET then computes the energy consumption $C(t)$ using:  $C(t)=\int_{0}^{t_\textit{simulation}} (r_m \circ c) (t) + (t_m \circ c) (t) \textit{dt}$.
%
The function $m$ is given and INET computes both $r_m$ and $t_m$ during the simulation.}

\smallskip
\noindent
\textbf{Algorithm-specific parameters.}
%\subsection*{Algorithm-specific parameters}
The behavior of each protocol depends on its parameters, for which we favor values previously reported in the literature.
%
For instance, we use 0.5s as the \emph{hello time} for \emph{MPR},  \emph{ProbFlood} and \emph{CDS-based}.
%
In the case of \emph{ABBA}, we use 0.3s as the maximum waiting time.
%
Finally, for \emph{ProbFlood} we use $k=15$ and $\sigma=1$.

\smallskip
\noindent
\textbf{Simulation time.}
%\subsection*{Simulation time}
As mentioned in Section~\ref{sec:background}, some algorithms first need to build a topology.
%
To guarantee that the first broadcast message is properly handled by all protocols, we set an initial warm-up period of 5s, allowing CDS-based algorithms to perform this initial construction.
%
%This time is used by  protocols can use to create any initial structure.
\deleted{
Likewise, \deleted{in our evaluation,} we stop each experiment a fixed amount of time (5s) after the beginning of the last broadcast operation.
%
This grace period allows the last message to reach all nodes.
}
%
The total simulation time is then $10 + {\delta t} \times M$, where ${\delta t}$ is the interval between broadcast messages and $M$ is the number of scheduled messages.


%\raziel{This definition (density) must be mentioned at the Related Work (or maybe Introduction), here it is too late}




 
%\deleted{


%\subsection*{Other such as battery and power consumption of the radio}


%\begin{itemize}
%	\item Experimental Setup
%	\begin{itemize}
%		\item Reproducibility of the experiments
%		\item Fairness of the experiments
%		\item Number of executions
%	\end{itemize}
%\end{itemize}
%}
